\chapter{Delivered MVP Alignment Notes}

One of the easiest ways to lose marks on this assignment is to submit a report and a codebase that look
like they were built independently. The final implementation therefore needs to be explained in relation
to the original design assumptions rather than pretending the first draft was already identical to the
final MVP.

\section{Design Assumptions vs Delivered MVP}

\begin{table}[H]
\centering
\caption{Key alignment notes between target architecture and delivered MVP}
\label{tab:mvp-alignment}
\begin{tabularx}{\textwidth}{p{4.1cm}p{4.8cm}X}
\toprule
\textbf{Design Assumption} & \textbf{Delivered MVP} & \textbf{Reason for Difference} \\
\midrule
External identity provider (OAuth/OIDC style) & Local bearer-token login via \texttt{POST /auth/login} & Keeps the proof of concept feasible while preserving the same trust boundary and role model \\
General relational database target & SQLite shared by backend and realtime & Portable, deterministic local setup; easier demo reset, lower infrastructure risk \\
Potentially hosted LLM provider & Configurable stub or local LM Studio provider & Improves reliability, cost control, and test determinism without changing the client-facing contract \\
PoC limited to basic frontend-backend communication & Delivered MVP includes live sync, autosave, permissions, AI, versions, revert, export, Docker, and CI & Team continued implementation beyond the minimum requirement once the architecture stabilized \\
Abstract collaboration transport & Concrete Yjs + WebSocket realtime service & Needed to validate simultaneous editing feasibility and collaboration correctness early \\
Generic AI suggestion concept & Side-panel review with full apply, partial apply, reject, and version snapshots & Better product integration and closer alignment with the user-story requirements \\
\bottomrule
\end{tabularx}
\end{table}

\section{Why These Differences Are Defensible}

None of these differences are accidental. They are the result of architectural trade-offs made under the
assignment's constraints. For example, using SQLite instead of a larger operational database deployment
does not weaken the architecture's modularity; it lowers setup friction so the team can spend more effort
on requirements, contracts, collaboration correctness, and testability. Likewise, using local bearer-token
authentication instead of full OAuth/OIDC is acceptable in the MVP because the report still explains the
role model, trust boundary, and future identity-provider integration direction.

This section exists to make those trade-offs explicit. The architecture remains coherent because the same
modules, APIs, and responsibilities described in Part 2 are actually visible in the repository and the
demoable system, even when the implementation chooses a coursework-friendly version of a more complex
target technology.
